
\subsubsection{Research Questions}

Our experimental design addresses three research questions. RQ1 (Existence \& Convergence): Can joint DPO + attribute training converge without catastrophic objective interference? We hypothesize gradient angles will remain below 120°. RQ2 (Representation Encoding): Do shared representations encode task-relevant information for both objectives? We hypothesize $\geq 70$\% linear probing accuracy on preference classification. RQ3 (Bidirectional Performance): Can a single jointly trained model achieve meaningful performance on both dimensions simultaneously? We hypothesize the model will exceed proof-of-concept thresholds ($\geq 50$\% preference win rate, $\geq 60$\% attribute steering accuracy).

\subsubsection{Hypotheses}

We employ a two-hypothesis validation strategy: H-E1 (Existence \& Convergence) tests whether joint training is implementable and convergent at proof-of-concept scale through four gate criteria: (1) monotonic decrease in both losses, (2) preference win rate $\geq 50$\%, (3) attribute steering accuracy $\geq 60$\%, and (4) gradient angle <120°. H-M1 (Shared Representation Learning) tests whether the joint model learns representations encoding preference and attribute information through linear probing analysis on layer 47 representations (preference classification accuracy $\geq 70$\%).

\subsubsection{Datasets and Baselines}

Both HH-RLHF (161,000 preference pairs) and OpenAssistant (88,000 attribute annotations) are publicly available via HuggingFace Datasets and were successfully verified during implementation. Our proof-of-concept experiments reference standalone baseline performance from prior work: DPO standalone achieves 57.5\% win rate \citep{Rafailov2023DPO}, representing the upper bound for preference alignment without attribute conditioning. SteerLM standalone achieves 87\% steering accuracy \citep{Dong2023SteerLM}, establishing the upper bound for attribute control. A sequential training baseline (DPO followed by attribute fine-tuning) was not trained in this study.

\subsubsection{Implementation}

Training configuration used GPT-2 XL with L\_total = 0.$7\cdot L_DPO$ + 0.$3\cdot L_attr$ loss formulation, AdamW optimizer at learning rate 1e-5, batch size 4 per GPU, and maximum sequence length 256 tokens over 100 proof-of-concept steps (versus 15,000 in full specification). This proof-of-concept scale prioritizes rapid feasibility validation over performance optimization---loss curves show continued decrease at training termination, suggesting models had not fully converged.
